%% kapa: exam=1 round=28 year=2017 period=1 form=A subject=economics subject_ko=경제학원론 seq=27 number=67
%% source: Q-Net 감정평가사 기출문제 (cst00309) artlSeq=5199603 file=extracted/1008/5199603/2017년 제28회 감정평가사 1차시험 문제지/1교시 A형.pdf
%% conversion: kordoc 4.24.51 -> tools/gen1.py (auto); review=proofread-sonnet; flags=none
\begin{question}{67}
효율성 임금(efficiency wage) 이론에 따르면 기업은 노동자에게 균형임금보다 높은 수준의 임금을 지급한다. 옳은 것을 모두 고른 것은?
\begin{framedtext}
ㄱ. 노동자의 생산성을 높일 수 있다. \par
ㄴ. 노동자의 근무태만이 늘어난다. \par
ㄷ. 노동자의 이직률을 낮출 수 있다.
\end{framedtext}
\begin{choices}
  \item ㄷ
  \item ㄱ, ㄴ
  \item ㄱ, ㄷ
  \item ㄴ, ㄷ
  \item ㄱ, ㄴ, ㄷ
\end{choices}
\end{question}
\answer{3}
